\section{Major formalization nodes}

The labels M01--M10 and their Lean facade modules are stable.  The checklists
inside each node are provisional implementation guides rather than additional
global dependency nodes.  Major nodes receive a formalization-complete marker
only after the Lean declarations and Blueprint links satisfy the project green
policy.

\begin{definition}[\nodeid{M01}: Finite-dimensional inertia reduction]
\label{node:m01}
\notready
\lean{ZetaZero.FiniteDimensionalInertiaReduction.rank_plus_hilbertSchmidt_minmax_after_linearEquiv}
The Lean core now formalizes codimension restriction, finite-rank deletion,
Hermitian spectral thresholding, the Hilbert--Schmidt dimension budget, and
transport of the resulting min--max theorem back to physical coordinates.
Blueprint statement/proof certification is the remaining integration step, so
this node is not yet marked green.  Paper source:
\texttt{sections/13\_minmax.tex}.
\end{definition}

\begin{definition}[\nodeid{M02}: Hardy gauge and invariant contour form]
\label{node:m02}
\notready
\lean{ZetaZero.HardyGaugeInvariantContourForm.deriv_hardyGauge_eq_branch_mul_zetaOne}
The analytic Hardy-gauge layer is now formalized: the functional-equation
factor, holomorphic square-root branch on an open high rectangle, exact
reflection law, and derivative identity $\mathscr Z'=qZ_1$.  The symmetric-shift
curvature identity is also formalized at the two-jet algebra level.  The
analytic second-shift derivative and invariant contour-integral representation
remain before M02 can be marked green.  Paper source: \texttt{sec:common-matrix}.
\end{definition}

\begin{definition}[\nodeid{M03}: Zero-side stationary geometry and packet evaluation]
\label{node:m03}
\notready
\uses{node:m02}
\lean{ZetaZero.Analytic.offCritical_reflection_pair}
The zero-side analytic foundation now formalizes zeta-zero multiplicity and
finite ordinate windows, critical-line reflection with multiplicity-preserving
off-line pairs, and the finite-set avoidance lemma used to choose good
horizontal heights.  Stationary residue blocks, the $Z_1$ local Jensen and
logarithmic-derivative estimates, exact packet evaluation rank, packet Riesz
bounds, and zero-side inertia bookkeeping remain.  Paper source:
\texttt{sec:zero-side}.
\end{definition}

\begin{definition}[\nodeid{M04}: Log-only right-edge source and stationary dilation]
\label{node:m04}
\notready
\uses{node:m02}
Formalize the Hermitian reflected source, closed-contour elimination of the
archimedean self block, the exact split
$L_1=H_{\log}+\Delta H_{\rm dir}+f'/(f-P)$, the
$O(L^{-2})$ direct residual, the physical one-character stationary carrier,
and exact dilation covariance.  Paper sources: \texttt{sec:right-edge} and
\texttt{sec:stationary}.
\end{definition}

\begin{definition}[\nodeid{M05}: Scalar HLZ and lagwise log-only principal operator]
\label{node:m05}
\notready
\uses{node:m04}
\lean{ZetaZero.HLPLocalModel.mangoldtConvolutionPower_support}
The existing Lean arithmetic library still contains the earlier von Mangoldt
convolution hierarchy, but the active manuscript target has changed.  The M05
interface to be formalized is now the separated-torus scalar three-$Z$ identity,
the uniform bounded-support Mellin-weight diagonal, local-height rescaling, and
the packet lift at the actual continuous lag $d=y-x$ with cutoff
$\ell\le Xe^{-|d|}$.  The target output is the two-face log-only kernel, not a
one-sided common-output multiplier.  Paper sources:
\texttt{sections/06\_scalar\_hlz\_core.tex} and
\texttt{sections/06\_lagwise\_packet\_hlz.tex}.
\end{definition}

\begin{definition}[\nodeid{M06}: Log-only global model and spectral floor]
\label{node:m06}
\notready
\uses{node:m05}
Formalize the corrected closed profile $G_{\log}$, including its
$\sinh(2\rho)$ term, the primitive identity, the rational Schur certificate,
the concrete lower bound, compression, and the conservative interface
$c_0=0.02$ used by the endgame.  Paper source:
\texttt{sections/08\_global\_model\_floor.tex}.
\end{definition}

\begin{definition}[\nodeid{M07}: Auxiliary divisor-Gram library]
\label{node:m07}
\notready
The existing Lean gcd/divisor Gram and arithmetic-transfer lemmas are retained
as an auxiliary library, but the revised invariant-source manuscript does not
route the main theorem through an outer divisor deformation.  M07 is therefore
not a dependency of the active M08 theorem path.  Its concrete arithmetic
formalization may continue independently.
\end{definition}

\begin{definition}[\nodeid{M08}: Geometric shell frame and primitive remainder]
\label{node:m08}
\notready
\uses{node:m03,node:m04,node:m06}
Formalize the stable-core frame, logarithmic unitary, global/local primitive
comparison, geometric exterior/terminal leakage, the two-variable
absolute-lag primitive criterion, and the log-only remainder classes.  No
macroscopic compressed-shift principal operator or HLP exact--model residual is
part of this node.  Paper sources:
\texttt{sections/10\_geometric\_shell\_core.tex},
\texttt{sections/11\_primitive\_remainder\_core.tex}, and
\texttt{sections/11\_logonly\_remainder.tex}.
\end{definition}

\begin{definition}[\nodeid{M09}: Retained space and parameter hierarchy]
\label{node:m09}
\notready
\uses{node:m08}
Formalize the reduced spectral space, global/local mean constraints,
accumulated codimension bounds, and the simplified acyclic parameter hierarchy.
Paper source: \texttt{sections/12\_retained\_space.tex}.
\end{definition}

\begin{definition}[\nodeid{M10}: Generic perturbation, min--max, and density-one endgame]
\label{node:m10}
\notready
\uses{node:m01,node:m03,node:m09}
\lean{ZetaZero.GenericPerturbationMinMaxEndgame.density_one_critical_line}
Formalize multiplicity-removing perturbation, the final finite-dimensional
min--max step, density-one selection, and the main theorem.  Paper sources:
\texttt{sections/00\_main\_theorem.tex},
\texttt{sections/13\_minmax.tex}, and
\texttt{sections/14\_endgame.tex}.
\end{definition}
